\chapter{Generative histories in the UNNS Substrate}\label{ch:intro}

\section{From values to generated occurrences}
The UNNS (Unbounded Nested Number Sequences) Substrate treats a number not only as a value but as a \emph{generated occurrence}: something that arises when admissible operators are applied in sequence, and that carries the record of how it arose. The operator studied here, SELECT-II, is
\[
T_N(m)\;=\;m\,\frac{(N+1)^2}{N},\qquad \text{defined when } N\mid m ,
\]
so that a \emph{history} is a source $m_0$ together with a word of labels $N_1,\dots,N_n$ for which every intermediate state is an integer divisible by the next label, and its \emph{endpoint} is $m_0\prod_i (N_i+1)^2/N_i$. The endpoint $72$ already shows what is at stake. It has exactly four histories,
\[
(1;1,1,2),\quad(1;1,2,1),\quad(3;3,2),\quad(10;5),
\]
two of which differ only by exchanging the adjacent labels $1$ and $2$, while the other two start from different sources. A reader who sees only the endpoint $72$ cannot tell which of the four produced it. The monograph asks how much of the history survives in the endpoint and in the intermediate observations one is willing to make, and what mathematical structure answers that question.

Four notions are used throughout.
\begin{itemize}[nosep]
\item The \emph{fibre} $F_t$ of an endpoint $t$ is the set of its histories, and $A(t)=\log_2|F_t|$ is its \emph{ancestral ambiguity}.
\item For a fixed source and a fixed set of $k$ distinct labels, the histories are \emph{orderings} of the labels, i.e.\ vertices of the permutohedron $P_k$; which orderings are \emph{admissible} depends on the divisibility of the source.
\item The admissible orderings span a subcomplex $K_A$ of the faces of $P_k$, the \emph{history complex}, whose homology records loops of histories that cannot be shrunk by exchanging neighbouring independent steps, hollow shells, and their higher analogues.
\item The orderings \emph{not} admitted form the \emph{forbidden complex} $Y_A$, which turns out to be the natural place in which to read the topology of $K_A$.
\end{itemize}

\section{What was found}\label{sec:found}
The results are collected in \cref{tab:found}; each row names the mechanism and where it is proved. The abstract states the same findings in prose. The development has a single thread: arithmetic admissibility is \emph{prefix-defined}, and everything else is a consequence of that fact together with the geometry of the permutohedron.
\begin{table}[ht]\centering\small
\begin{tabularx}{\textwidth}{@{}>{\raggedright\arraybackslash}p{3.0cm}>{\raggedright\arraybackslash}X>{\raggedright\arraybackslash}p{1.9cm}@{}}
\toprule
Discovery & Statement & Proved in\\\midrule
Prefix law & A history is admissible iff the denominator of the product of $(N{+}1)^2/N$ over every prefix set divides the source. Admissible sets form an lcm-closure system with a unique minimal source, and the history complexes grow monotonically along divisibility. & \cref{ch:arith}\\
Reconstruction & The generator-coloured swap graph determines the whole history complex. The uncoloured graph determines its two-skeleton for prefix-defined systems. & \cref{ch:swaps}\\
Arithmetic activation of holes & An explicit 96-face two-cycle survives exactly when its own prefix sets are admissible and none of four minimal families of further prefix sets is entirely admissible. For the anchor contents this is an explicit gcd condition on the source; a strong form of gate necessity is false. & \cref{ch:topology}\\
Survival law of every cycle & For any cycle of any permutohedron, survival is a monotone Boolean function of the admissible prefix sets; the atlas of such laws organises the two-dimensional holes of a six-label census. & \cref{ch:atlas}\\
Lifting & A suspension-type calculus raises holes from $k$ to $k{+}1$ operators; diagonal lifts never carry holes; the dependence of a lift on the filling is governed by an exact criterion; certified holes occur in dimensions three and four. & \cref{ch:dimension}\\
Information layers & $A(t)$ decomposes exactly into source, content, class and within-class parts; an exact radius-two identity bounds what local statistics see; equal counts, entropies and local statistics through radius two are compatible with different Betti numbers. & \cref{ch:ancestry,ch:access,ch:separation}\\
Forbidden-history duality & $\beta_2(K_A)=\dim H_1(Y_A)=\dim H_1(N[D_A])$, with $N[D_A]$ the nerve of the dead prefix sets; a strong-collapse core decides, count-level statistics do not. & \cref{ch:duality,ch:minimal}\\
Locality & For six events, constraints on consecutive pairs cannot create a two-dimensional hole; constraints on consecutive triples can create five. & \cref{ch:locality}\\
\bottomrule
\end{tabularx}
\caption{The principal discoveries and the chapters in which they are proved.}\label{tab:found}
\end{table}

\section{Two views of one system}
The monograph moves between two descriptions of the same system, shown in \cref{fig:overview}. On the \emph{generation} side, a history is successively projected to its swap class, to its source and content, and to its endpoint; each projection forgets something, and \cref{ch:ancestry} measures exactly what. On the \emph{geometry} side, the admissible orderings form a subcomplex of the permutohedron; Alexander duality turns it into the forbidden complex, the nerve theorem into a small complex on the dead prefix sets, and strong collapse into its core. The SELECT-II rule is the bridge: it fixes the admissible set from the source through the prefix law, and divisibility of the source moves the system along the geometry side.

\begin{figure}[ht]\centering
\begin{tikzpicture}[>=Stealth,font=\small,
  bx/.style={draw=ink,rounded corners=2pt,align=center,minimum height=8.5mm,fill=white,inner sep=3pt},
  hl/.style={draw=accent,thick,rounded corners=2pt,align=center,minimum height=8.5mm,fill=accent!8,inner sep=3pt},
  lb/.style={font=\scriptsize\itshape,text=ink!75,align=center}]
\node[lb] at (-4.4,3.6) {\normalsize\upshape\bfseries Generation (ancestry projections)};
\node[bx] (W) at (-4.4,2.4) {history $W$\\[-1pt]{\scriptsize source, word}};
\node[bx] (C) at (-4.4,1.0) {swap class $C$};
\node[bx] (MQ) at (-4.4,-0.4) {source and content $(M,Q)$};
\node[bx] (T) at (-4.4,-1.8) {endpoint $T$};
\draw[->] (W)--node[right,lb]{forget order inside a class}(C);
\draw[->] (C)--node[right,lb]{forget the class}(MQ);
\draw[->] (MQ)--node[right,lb]{multiply out}(T);
\node[lb] at (4.0,3.6) {\normalsize\upshape\bfseries Geometry (admissibility)};
\node[hl] (A) at (4.0,2.4) {admissible set $A\subseteq S_k$};
\node[bx] (K) at (4.0,1.0) {history complex $K_A\subseteq\partial P_k$};
\node[bx] (Y) at (4.0,-0.4) {forbidden complex $Y_A$};
\node[bx] (N) at (4.0,-1.8) {dead-set nerve $N[D_A]$, its core};
\draw[->] (A)--node[right,lb]{faces inside $A$}(K);
\draw[<->] (K)--node[right,lb]{Alexander duality}(Y);
\draw[->] (Y)--node[right,lb]{nerve theorem}(N);
\node[hl] (R) at (0,2.4) {SELECT-II};
\draw[->,accent,thick] (R.east)--(A.west);
\draw[->,accent,thick] (R.west)--(W.east);
\end{tikzpicture}
\caption{The two views of a system. The left column lists what is forgotten when a history is observed only through coarser data; the right column lists the structures through which its topology is read. The operator SELECT-II, through its prefix law and the divisibility of the source, determines the admissible set. Observation levels and what each recovers are developed in \cref{ch:access,ch:separation}.}\label{fig:overview}
\end{figure}

\section{Organisation}
\Cref{ch:grammars} fixes the setting and states which constructions need arithmetic and which do not. \Cref{ch:ancestry} treats the multiplicity of histories, \cref{ch:arith} the arithmetic that selects admissible orderings, and \cref{ch:swaps} the geometry of legal swaps. \Cref{ch:topology,ch:atlas,ch:dimension} study holes, their survival laws and their lifts. \Cref{ch:access,ch:separation} ask which observations recover which of these, \cref{ch:duality,ch:minimal} read the topology from what is forbidden, and \cref{ch:locality} separates the contribution of arithmetic from that of locality. \Cref{ch:synthesis} collects the principal results, their relation to established mathematics and the implications for the Substrate. A reader interested in the duality alone can read \cref{ch:grammars}, the definitions of \cref{ch:swaps}, and then \cref{ch:duality,ch:minimal,ch:locality}.

\section{Conventions}
Homology and cohomology have coefficients in $\F=\mathrm{GF}(2)$ unless stated. All computations are exact: integers, rationals and linear algebra over $\mathrm{GF}(2)$. Where a result rests on a finite computation, or holds only on a stated finite domain, its title says so (\emph{computer-assisted}, \emph{finite-exact}, \emph{bounded}, \emph{falsified}); \cref{app:status} defines the tags and lists what was re-derived independently for this monograph. Observations about finite censuses are never stated as theorems. The programme-internal designations of the source releases (GIG-I, GIG-INFO, SELECT-II) occur only here, in the notation of \cref{ch:grammars} and in \cref{app:status}, where they identify sources.
